% ============================================
% Johns Hopkins Letter Template
% Assistant Professor Cover Letter – Georgia State University
% ============================================

\documentclass[11pt,letterpaper]{letter}

\usepackage{graphicx}
\usepackage{eso-pic}
\usepackage{geometry}
\usepackage{microtype}
\usepackage[hidelinks]{hyperref}

% ----- Page Setup -----
\geometry{left=25mm,right=25mm,top=25mm,bottom=25mm}

% ----- Logo Placement (first page only) -----
\newcommand{\PutJHULogoFG}[1]{%
  \AddToShipoutPictureFG*{%
    \AtPageUpperLeft{%
      \hspace{10mm}%
      \raisebox{-38mm}{%
        \includegraphics[width=6cm]{#1}%
      }%
    }%
  }%
}

% ----- Sender Information -----
\signature{Decory Edwards}
\address{Department of Economics \\ Johns Hopkins University \\ Baltimore, MD 21218 \\ \texttt{dedwar65@jh.edu}}

\begin{document}
\PutJHULogoFG{Images/jhulogo.png}

\begin{letter}{Search Committee \\
Department of Economics \\
Andrew Young School of Policy Studies \\
Georgia State University \\
Atlanta, GA 30303 \\ USA}

\opening{Dear Members of the Search Committee,}

I am writing to express my interest in the tenure-track Assistant Professor position in Econometrics in the Department of Economics at the Andrew Young School of Policy Studies at Georgia State University, beginning August 2026. I am a Ph.D. candidate in Economics at Johns Hopkins University, advised by Professors Christopher D.~Carroll, Nicholas W.~Papageorge, and Stelios Fourakis, and I expect to receive my degree in May 2026. My research fields are macroeconomics, wealth and income dynamics, and computational economics.

My research agenda focuses on understanding the determinants of wealth inequality and the mechanisms through which heterogeneity in returns to wealth shapes aggregate outcomes. In my job-market paper, I estimate the degree of heterogeneity in individual portfolio returns using micro-data from the Survey of Consumer Finances and simulate a structural model in which households face idiosyncratic return risk. I show that allowing for persistent heterogeneity in returns substantially improves the model’s ability to match the empirical distribution of wealth, offering a tractable way to connect micro-level return dispersion to macroeconomic inequality.

In a second project, I use the Health and Retirement Study to examine the relationship between interpersonal trust and portfolio performance. Building on the literature that finds a hump-shaped relationship between trust and income, I show that a similar relationship holds for individual returns to wealth measured in the HRS data, highlighting a behavioral channel through which social attitudes shape saving and investment outcomes. This work also provides new estimates of the persistent component of returns to net worth, which motivate the calibration of heterogeneity in my structural model.

Georgia State University’s Economics Department—internationally recognized for applied microeconometrics, policy analysis, and interdisciplinary research—provides an ideal setting for my quantitative, data-oriented research. I am particularly drawn to the position’s emphasis on graduate instruction in the Econometrics sequence and the school’s strengths in applied microeconomics, causal inference, and policy evaluation. My background in computational modeling, empirical macroeconomics, and large-scale microdata aligns closely with GSU’s mission to advance research on health, labor, public finance, and urban economics. I would welcome opportunities to collaborate with the Georgia Policy Labs, the Georgia Health Policy Center, and researchers at the Federal Reserve Bank of Atlanta.

\textbf{Teaching.} My teaching experience centers on undergraduate macroeconomics and an upper-level macro policy seminar, where I have led weekly discussion sections and held regular office hours for classes ranging from 16 to 40 students. My approach is student-centered: I aim to help students “convince themselves” that they have moved from confusion to understanding by holding structured office-hour coaching that builds trust and confidence. At Georgia State, I am prepared to teach graduate and undergraduate \textit{Econometrics}, applied microeconometrics, and courses focused on machine learning and data-intensive empirical methods, as well as contribute to the broader curriculum in economics. I also welcome the opportunity to mentor graduate students and support the department’s experiential and policy-engaged learning initiatives.

I believe I would be a strong fit because my empirical and computational training allows me to (i) integrate modern econometric tools and machine learning techniques into applied economic research, (ii) provide rigorous, hands-on instruction using real-world datasets, and (iii) contribute to GSU’s mission as an R1 university committed to research excellence, student mobility, and serving a diverse and global student body. I am enthusiastic about contributing to the Andrew Young School’s collaborative research environment and its commitment to policy-relevant scholarship.

I would be honored to join the Department of Economics at Georgia State University. Thank you for your consideration; I welcome the opportunity to discuss my fit with your department at your convenience.

\closing{Sincerely,}

\end{letter}
\end{document}
